\documentclass{article}
\usepackage[utf8]{inputenc}
\usepackage{graphicx,grffile}


\title{ARG2 procedure}
\author{ Andreas }
\date{November 2019}

\begin{document}
\maketitle



\section{Summary}

When deciding how to fill vacancies in a team with different roles, it is important to understand what skills are needed for each role and how to identify these skills in a set of candidates.
Efficient search for information on candidates' skills means being adaptive and focusing information queries on relevant instead of irrelevant features.
This study aims to investigate when this adaptiveness in information search emerges and how it develops across childhood.
In a game setup, children are presented with a set of cartoon monsters that perform 2 different sports activities (throwing and jumping), and with 2 body features that enable them for either sport (long arms and long legs).
Children learn how each body feature maps to ability in one sport, and how these features are distributed in the monster population---either evenly (Uniform condition), or unevenly (Skewed condition).
Children's task is to fill two slots in a team with a ``thrower" and a ``jumper" by finding monsters with the relevant features.
For that, we draw monsters from the population one by one with arms and legs covered, and children can disclose these features at a cost that they pay from an endowed amount of stickers or marbles. \\
\textbf{DV}: We measure the mean informativeness of children's queries into monsters' features and the accuracy of their final choices of monsters to fill the slots. \\
\textbf{IV}: Independent variables are age (range: 4 -- 10 years) and condition (within subjects, varied across 2 rounds of the game). \\
We counterbalance which body feature is rare and how monsters look across conditions.



\section{Design}\label{design}

\subsection{Independent variables:}\label{independent-variables}

\begin{itemize}

\item
  \textbf{age}: continuous, range: 4;0 \ldots{} 8;0
\item
  \textbf{condition}: uniform vs.~skewed (between subjects)
\end{itemize}

\subsection{Counterbalancing:}\label{counterbalancing}

\begin{itemize}

\item
  order
\item
  rare bodypart
\end{itemize}

\section{Materials}\label{materials}

\begin{itemize}

\item
  playing cards (quantity to be determined) with arms and legs covered (1)
\item
  demo cards: 1 long-arm, 1 long-leg, 1 short-arm, 1 short-leg monster, serve as demonstration monsters at the demo stage; at start, limbs covered (2)
\item
  ``trash can'': small paper basket/bucket for disused monsters (3)
\item
  ``target'' background picture with thrower/jumper slot marked by bucket/hurdle (4)
\item
  marbles or stickers (5)
\item
  2 bowls as containers for marbles/stickers (6)
\item
  laptop with ``card reader'' for the playing cards (7)
\item
  camera* (8)
\end{itemize}

*) It may be better to cover the laptop camera and use an overt camera to not give the impression that we are trying to hide our video recording.
This may be important with some parents.

\begin{figure}
\centering
\includegraphics[width=.5\textwidth]{images/setup.png}
\caption{ The setup (bird's eye view). Numbers are referenced in the Materials section. C: Child; E: Experimenter. }
\label{fig:setup}
\end{figure}



\section{Familiarization}\label{familiarization}

\begin{figure}
\centering
\includegraphics[width=.8\textwidth]{images/card_construction.png}
\caption{ How we could construct playing cards with hideable features. Purple and green background colors are only for illustration, they will be equal on the material. }
\label{fig:cards}
\end{figure}

Experimenter (E) and child (C) sit down side by side in front of the table, with the laptop and case with playing materials on it (see fig. \ref{fig:setup}).
\emph{``Guck mal, wir haben hier nämlich ein cooles Spiel mitgebracht.''}
E picks up the playing cards, scatters them on the table, with monsters' limbs open.
\emph{``Siehst du die ganzen Monster hier?
Die kommen alle vom Planeten Seeblick.
Die sind alle total sportlich und die haben zwei Lieblingssportarten.''} \\

Counterbalance: \{\textbf{jumping first} \textbar{} throwing first\}



\subsection{Demo 1: jumping, "Nachgucken"}\label{demo-1-jumping-nachgucken}

\begin{itemize}
\item
  \emph{``Das hier ist der eine Sport.''} Click to first image on laptop (hurdle). \emph{``Die machen voll gerne Hindernisspringen. Dabei versuchen die, es über die Hürde hier {[}point{]} zu schaffen. Und da gibts Monster, die schaffen das ganz gut, und Monster, die schaffen das garnicht. Und wir wollen rausfinden, welche es gut schaffen.''}
\item
  E picks up the two jumping demo monsters, with features hidden. \emph{``Nehmen wir mal die beiden hier. Guck mal, bei denen sind noch die Arme und Beine versteckt. Wollen wir da mal nachgucken?''}
\item
  E opens all 4 limbs, comments \emph{``Da sind die Arme, {[}\ldots{}{]} da sind die Beine {[}\ldots{}{]} So und denen gucken wir mal dabei zu wie die über ein Hindernis springen, ja?''}
\item
  E puts the short-leg monster card into the card reader, clicks to the first jumping video, in which the monster fails. \emph{``Oh nein, das Monster hats leider nicht übers Hindernis geschafft, weil das hat ja auch viel zu kurze Beine dafür! Siehst du das? {[}point{]} Monster mit kurzen Beinen sind immer ganz schlecht beim Hindernisspringen.''}
\item
  \emph{``Jetzt gucken wir mal wie sich das Monster hier anstellt.''} E puts the long-leg monster card into the card reader, clicks to the second jumping video, in which the monster succeeds. \emph{``Yay, das andere Monster hats geschafft! Weil das hat ja auch ganz lange Beine, siehst du? {[}point{]} Monster mit langen Beinen sind immer total gut beim Hindernisspringen.''}
\item
  E puts both demo cards side by side. \emph{``Guck mal, hier sind die beiden nochmal. Komm wir spielen jetzt mal ein Spiel. Du suchst dir eins von den beiden Monstern aus, und wenn es über das Hindernis kommt, dann kriegst du 5 Murmeln. Und die Murmeln, die brauchst du nämlich gleich noch in unserem Spiel.''} C should select the long-leg monster; if not, E convinces C of the long-leg monster. \textbf{Note that this can be a drop criterion.}
\item
  E puts the long-leg monster into the card reader, clicks to the next video, both watch together. \emph{``Super gemacht! Dafür gewinnst du jetzt 5 Murmeln.''} E puts 5 marbles into C's bowl.
\item
  E puts both demo monsters into the trash can, comments \emph{``So, die beiden hatten jetzt ihren Spaß.''}
\end{itemize}



\subsection{Demo 2: Throwing, buying info, repeat ``Nachgucken"}\label{demo-2-throwing-buying-info-repeat-nachgucken}

\begin{itemize}
\item
  \emph{``Weißt du welchen anderen Sport die Monster noch total gerne machen? Guck mal hier.''} Click to second image on laptop (bucket \& ball). \emph{``Die machen voll gerne Zielwerfen. Dabei versuchen die, den Ball von ganz weit weg in den Eimer {[}point{]} zu werfen. Und dabei gibts Monster, die schaffen das ganz gut, und Monster, die schaffen das garnicht. Und wir wollen wieder rausfinden, welche es gut schaffen.''}
\item
  E picks up the two throwing demo monsters. \emph{``Nehmen wir mal die beiden hier. Da sind ja auch noch die Arme und Beine versteckt. Wollen wir da mal nachgucken? {[}wait for response{]} Ok, aber weißt du was, ab jetzt haben wir eine Spielregel. Wenn du willst, dann kannst du nachgucken was das Monster für Arme hat, aber dafür bekomme ich eine Murmel von dir. Wollen wir bei \textbf{dem} {[}short{]} Monster mal nachgucken was das für Arme hat?''}
\item
  E awaits C's agreement. \emph{``Okay, dann bekomm ich jetzt eine Murmel von dir bitte.''}
\item
  E opens the arms. \emph{``Da sind die Arme vom einen Monster. Dann gucken wir dem mal beim Werfen zu.''} E puts the selected monster card into the card reader, clicks to the first throwing video, in which the monster fails. \emph{``Oh nein, das Monster hat den Eimer leider nicht getroffen, weil das hat ja auch viel zu kurze Arme dafür! Siehst du das? {[}point{]} Monster mit kurzen Armen sind immer ganz schlecht beim Zielwerfen.''}
\item
  \emph{``Los, wir müssen wohl mal gucken was das andere Monster für Arme hat, ja?''} E awaits C's agreement. \emph{``Okay, dann bekomm ich eine Murmel von dir bitte.''}
\item
  E opens the arms. \emph{``Da sind die Arme vom zweiten Monster. Dann gucken wir doch mal.''} E puts the selected monster card into the card reader, clicks to the second throwing video, in which the monster succeeds. \emph{``Yay, das andere Monster hats geschafft! Weil das hat ja auch ganz lange Arme, siehst du? {[}point{]} Monster mit langen Armen sind immer total gut beim Zielwerfen.''}
\item
  E puts both demo cards side by side. \emph{``Guck mal, hier sind die beiden nochmal. Komm wir machen das Spiel nochmal. Du suchst dir eins von den beiden Monstern aus, und wenn es den Eimer trifft, dann bekommst du wieder 5 Murmeln in unserem Spiel.''} C should select the long-arm monster; if not, E convinces C of the long-arm monster. \textbf{Drop criterion.}
\item
  E puts the long-arm monster into the card reader, clicks to the next video, both watch together. \emph{``Wieder richtig gemacht! {[}high five{]} Da gewinnst du jetzt wieder 5 Murmeln.''} E puts 5 marbles into C's bowl.
\item
  E puts both demo monsters into the trash can, comments \emph{``So, die beiden hatten jetzt auch ihren Spaß.''}
\end{itemize}

\subsection{Study population}

\begin{itemize}
\item
  E focuses C's attention to the playing cards.
  \emph{``Jetzt kommt der spannende Teil.
  Hier haben wir eine ganze Menge von Monstern vom Planeten Seeblick.
  Und du siehst ja, manche haben kurze Arme, manche haben lange Arme.
  Und manche haben kurze Beine, manche haben lange Beine.
  Los wir machen mal Ordnung und tun alle mit langen Armen* hierhin [places first] auf einen Stapel und alle mit kurzen Armen daneben; machst du mit?"}
  If C does not cooperate, give C a card and ask them to place it on the right stack. \\
  {*}) In the Skewed condition, the non-skewed body part.
\item
  When done: \emph{``Siehst du das? Monster mit langen Beinen haben wir hier genauso viele wie Monster mit kurzen Beinen.
  Dann legen wir jetzt mal alle Monster mit langen Armen hierhin [places first] und alle Monster mit kurzen Armen daneben, los gehts!"}
  Again, facilitate cooperation if necessary. \\
  \textbf{Uniform condition}: Comment on equal distribution again. \\
  \textbf{Skewed condition}: Oh, siehst du das? Von den Monstern mit langen Beinen gibts ja nur zwei und alle anderen haben kurze Beine!
\end{itemize}



\subsection{Hide population, instruct}

\begin{itemize}
\item
  E starts to close all playing cards. \emph{So und jetzt machen wir die alle mal zu.}
\item
  When done closing cards, E ``shuffles'' all cards on the table, then stacks them and spreads them face down in a grid (so we can see how many are left), while saying:
  \emph{``Jetzt mischen wir mal alle schön durcheinander und tun die hier hin.''}
\item
  E instructs: \emph{``Und jetzt pass mal auf, ich erklär dir mal wie unser Spiel geht.''}
  E places the target picture in front of C.
\item
  \emph{``Du sollst heute eine Monstermannschaft zusammenstellen.
  Guck mal hier {[}point to target{]}, für deine Mannschaft, da brauchst du einen Werfer und einen Springer.
  Und du willst ja bestimmt zwei Monster für die Aufgaben finden, die das richtig gut können, stimmts?
  Denn wenn du eine richtig gute Mannschaft zusammenstellst, dann gewinnst du am Ende was tolles!''}
\item
  E draws 1 monster.
  \emph{``Als erstes ziehe ich ein Monster.
  Und du sollst rausfinden, ob das Monster ein guter Mitspieler in deiner Monstermannschaft ist, also ob du es dahin legen willst {[}point: targets{]}.
  Jetzt können wir ja garnicht sehen, was es für Arme und Beine hat, stimmts?
  Wenn du willst, dann kannst du nachgucken.
  Und wenn du willst, dann leg ich dir ein neues Monster aus dem Stapel hin und wir tun das alte weg.''} \\
  
\end{itemize}

\section{Test}\label{test}

\begin{itemize}
\item
  Begin: \emph{``Also, was möchtest du machen?
  Was bei dem Monster hier nachgucken?
  Oder das Monster wegtun und ein anderes angucken?
  Oder das Monster in deine Mannschaft nehmen?''}
\item
  Repeat these options after every move.
\item
  When C chooses to place a monster on the target, ask \emph{``Und was soll das Monster in deiner Mannschaft machen, werfen oder springen?''}
\end{itemize}

\end{document}

